\chapter{Glossary of Core Definitions}\label{ch:auto-007}
This appendix collects first-pass definitions used throughout Chapters 1--10 and the advanced-topics appendix.

\begin{description}
    \item[\textit{Parallelism}] Doing multiple useful operations concurrently to increase throughput.
    \item[\textit{Speculation}] Executing work before all required conditions are confirmed, with recovery if assumptions are wrong.
    \item[\textit{Prediction}] Estimating uncertain future behavior to guide speculative execution.
    \item[\textit{Instruction-Level Parallelism (ILP)}] Overlap of independent instructions within one thread.
    \item[\textit{Very Long Instruction Word (VLIW)}] Instruction encoding style that packs several compiler-scheduled ops into one instruction word.
    \item[\textit{Polycyclic Scheduling}] Early name for overlapping multiple loop iterations in one compact repeating schedule.
    \item[\textit{Software Pipelining}] Loop-scheduling technique that overlaps operations from different iterations and emits a prologue, steady-state kernel, and epilogue.
    \item[\textit{Thread-Level Parallelism (TLP)}] Concurrent execution of multiple threads on one or more cores.
    \item[\textit{Processing Element (PE)}] One independently executing processor, core, or other execution agent in a multiprocessor.
    \item[\textit{Digital Signal Processor (DSP)}] Processor class optimized for regular numeric workloads such as filtering, radio, and image-processing kernels.
    \item[\textit{SIMD}] One instruction stream operating on multiple data elements in lockstep.
    \item[\textit{MIMD}] Multiple instruction streams operating on multiple data streams independently.
    \item[\textit{Shared Memory}] Multiprocessor organization in which all participating processors conceptually operate in one common address space.
    \item[\textit{Message Passing}] Multiprocessor organization in which processors communicate by explicitly sending data between private memories.
    \item[\textit{Partitioned Global Address Space (PGAS)}] Memory model with a global address space in which some regions are much more expensive to access than others.
    \item[\textit{Accelerator}] Specialized processing engine, often attached to a host CPU through an external link and using distinct memory.
    \item[\textit{Bus}] Shared observable broadcast medium across which all attached agents can see the same transactions.
    \item[\textit{Point-to-Point Interconnect}] Link organization in which packets move between specific endpoints rather than across one shared observable medium.
    \item[\textit{Cache Coherence}] Property that keeps multiple cached copies of the same location mutually consistent enough to preserve shared-memory semantics.
    \item[\textit{Snooping Coherence}] Coherence method in which each cache monitors shared interconnect traffic and reacts when another processor requests a block it holds.
    \item[\textit{Directory-Based Coherence}] Coherence method in which a directory tracks which caches currently hold a block and coordinates invalidations or forwarding.
    \item[\textit{MSI}] Coherence protocol whose stable states are modified, shared, and invalid.
    \item[\textit{MESI}] Coherence protocol whose stable states are modified, exclusive, shared, and invalid.
    \item[\textit{MOESI}] Coherence protocol whose stable states are modified, owned, exclusive, shared, and invalid.
    \item[\textit{MOESIF}] Coherence protocol whose stable states are modified, owned, exclusive, shared, invalid, and forward.
    \item[\textit{GetS}] Coherence request for a shared readable copy of a block.
    \item[\textit{GetM}] Coherence request for exclusive writable ownership of a block.
    \item[\textit{WriteBack}] Propagation of a dirty cache block to lower memory.
    \item[\textit{Exclusive (E)}] Clean state in which exactly one cache holds the block and may later write it without first broadcasting an upgrade request.
    \item[\textit{Owned (O)}] Dirty shared state in which one cache remains responsible for the current value and may respond to later requests even though other caches hold clean shared copies.
    \item[\textit{Forward (F)}] Clean shared state designating one cache as the responder for later shared requests.
    \item[\textit{Hit Line}] Shared indication line used to tell a requester that some other cache already holds the block, thereby distinguishing \texttt{E} from \texttt{S} in a snooping protocol.
    \item[\textit{Intervene}] Snooper action in which a cache supplies data directly in response to another processor's request.
    \item[\textit{Abort}] Snooper response instructing a requester to wait or retry because ownership transfer is not yet complete.
    \item[\textit{Transient State}] Intermediate coherence state used while a block is waiting for a request or ownership transfer to complete.
    \item[\textit{Presence Bit Vector}] Directory field that records which caches or processors currently hold a copy of one block.
    \item[\textit{Limited Directory}] Directory organization that records only a bounded number of sharers explicitly and uses an overflow policy when that bound is exceeded.
    \item[\textit{Recall}] Directory command requiring a cache to return the current block value and surrender ownership.
    \item[\textit{Invalidate}] Directory command requiring a cache to discard its current copy of a block.
    \item[\textit{Ack}] Directory-style acknowledgment indicating that a required coherence action completed.
    \item[\textit{Nack}] Directory-style negative acknowledgment indicating that a request cannot yet be completed and should be retried later.
    \item[\textit{Sharing Pattern}] Characteristic way in which a block is accessed across processors, such as private, read-only, producer-consumer, migratory, or read-write shared.
    \item[\textit{False Sharing}] Performance problem in which unrelated variables placed in the same cache block force unnecessary invalidations or transfers because one of them is actively shared.
    \item[\textit{Coherence Miss}] Miss caused by coherence actions such as invalidation or ownership transfer rather than by compulsory, capacity, or conflict effects alone.
    \item[\textit{Inclusive Cache Hierarchy}] Hierarchy in which data present in an upper level is guaranteed to also be present in the next lower level.
    \item[\textit{Exclusive Cache Hierarchy}] Hierarchy in which adjacent levels are intended not to hold duplicate copies of the same block.
    \item[\textit{Non-Exclusive Cache Hierarchy}] Hierarchy that allows overlap between adjacent levels without requiring it for every block.
    \item[\textit{Critical Section}] Region of code that must not be executed concurrently by multiple threads if shared data is to remain correct.
    \item[\textit{Lock}] Synchronization primitive that grants exclusive access to a critical section.
    \item[\textit{Mutex}] Mutual-exclusion lock allowing at most one thread into a protected critical section at a time.
    \item[\textit{Barrier}] Synchronization point that forces all participating threads to wait until the last one arrives.
    \item[\textit{Spin Lock}] Lock implementation in which a waiting thread repeatedly polls a shared variable instead of blocking.
    \item[\textit{Atomic Read-Modify-Write}] Indivisible memory operation that reads a location, transforms it, and writes the new value with no visible intermediate state.
    \item[\textit{Test-and-Set}] Atomic read-modify-write primitive that writes a 1 and returns the old value.
    \item[\textit{Test-and-Test-and-Set}] Spin-lock strategy that first polls with ordinary reads and attempts an atomic acquisition only when the lock appears free.
    \item[\textit{Load-Link / Store-Conditional (LL/SC)}] Atomic primitive pair in which a load establishes a reservation and a later store succeeds only if no conflicting coherence event invalidated that reservation.
    \item[\textit{Fetch-and-Op}] Atomic primitive that returns the old value of a location while applying some operation to it.
    \item[\textit{Fetch-and-Add}] Fetch-and-op specialization in which the atomic update is addition.
    \item[\textit{Tree Lock}] Lock organization that reduces contention by arranging locks hierarchically so contenders meet in stages.
    \item[\textit{Tournament Lock}] Another name for a hierarchical tree lock in which contenders advance through multiple levels.
    \item[\textit{Sense-Reversing Barrier}] Barrier that uses a phase bit so one barrier instance is not confused with the next.
    \item[\textit{OpenMP}] Standard shared-memory parallel-programming interface based largely on compiler directives and runtime support.
    \item[\textit{MPI}] Message Passing Interface; standard API for distributed-memory or message-passing parallel programming.
    \item[\textit{Transactional Memory}] Synchronization model in which a critical section executes optimistically and commits only if no conflicting access occurred.
    \item[\textit{Hardware Transactional Memory}] Transactional-memory implementation that uses hardware structures such as caches to track speculative read and write sets.
    \item[\textit{Software Transactional Memory}] Transactional-memory implementation that tracks conflicts and retries through software metadata rather than specialized hardware.
    \item[\textit{Memory Consistency}] Rules governing the visibility and ordering of memory operations to different addresses across multiple processors.
    \item[\textit{Strict Consistency}] Memory model in which all processors observe memory operations in one global order that matches actual time.
    \item[\textit{Sequential Consistency}] Strong memory model in which execution appears as one total interleaving that preserves each processor's program order.
    \item[\textit{Weak Consistency}] Family of relaxed models that provide strong ordering mainly at explicit synchronization operations rather than at every ordinary load and store.
    \item[\textit{Acquire}] Synchronization action that prevents later memory accesses from moving before it.
    \item[\textit{Release}] Synchronization action that prevents earlier memory accesses from moving after it.
    \item[\textit{Store Fence}] Ordering operation that ensures earlier stores become visible before later memory actions proceed.
    \item[\textit{Load Fence}] Ordering operation that prevents later loads from bypassing earlier outstanding memory actions.
    \item[\textit{Memory Fence}] Explicit ordering operation that constrains permitted memory reordering around a program point.
    \item[\textit{Flit}] Flow-control digit; the unit of transfer controlled by the network's flow-control mechanism.
    \item[\textit{K-ary N-Dimensional Topology}] Direct-network family with $n$ dimensions and $k$ nodes in each dimension, for a total of $k^n$ nodes.
    \item[\textit{Direct Topology}] Network topology in which processors or nodes are connected directly to selected neighboring nodes.
    \item[\textit{Indirect Topology}] Network topology in which processors communicate through a separate switching fabric.
    \item[\textit{Hypercube}] Direct topology with $2^d$ nodes in which each node has degree $d$ and neighbors differ in one coordinate.
    \item[\textit{Mesh}] Direct topology in which nodes connect to neighbors along one or more dimensions without wraparound edge links.
    \item[\textit{Torus}] Direct topology in which each dimension wraps around so the network forms one ring per dimension.
    \item[\textit{Gray Code}] Binary numbering scheme in which adjacent code words differ in exactly one bit.
    \item[\textit{Bisection Bandwidth}] Minimum aggregate link bandwidth crossing any equal-halves cut of a network.
    \item[\textit{Crossbar}] Indirect switching fabric that can connect any input to any output.
    \item[\textit{Multi-Stage Interconnect}] Indirect network built from several stages of small switches rather than one full crossbar.
    \item[\textit{Perfect Shuffle}] Fixed wiring permutation commonly used between stages of an omega network.
    \item[\textit{Omega Network}] Multi-stage interconnect in which one destination bit guides the routing choice at each stage.
    \item[\textit{Routing}] Process of choosing the path a flit or packet takes through the interconnect.
    \item[\textit{Deterministic Routing}] Routing policy that always chooses the same path for a given source and destination.
    \item[\textit{Dimension-Order Routing}] Deterministic routing policy that traverses dimensions in a fixed order such as $x$ before $y$.
    \item[\textit{Adaptive Routing}] Routing policy that chooses among multiple next hops based on current congestion information.
    \item[\textit{Oblivious Routing}] Non-adaptive routing policy that ignores current load but may still use randomized path selection for balance.
    \item[\textit{Valiant's Algorithm}] Oblivious routing method that sends traffic first to a randomly chosen intermediate node and then to the final destination.
    \item[\textit{Flow Control}] Rules determining when data may advance through the network and what happens when downstream resources are full.
    \item[\textit{Circuit Switching}] Communication method that reserves an entire route before data transmission begins.
    \item[\textit{Packet Switching}] Communication method in which independently routed packets carry their own destination information.
    \item[\textit{Packet}] Message unit that carries routing information and may be divided into multiple flits for transport.
    \item[\textit{Back Pressure}] Upstream throttling mechanism used to prevent downstream buffers from overflowing.
    \item[\textit{Credit-Based Flow Control}] Back-pressure scheme in which a sender tracks how many downstream buffer slots remain available.
    \item[\textit{Store-and-Forward}] Flow-control method that buffers the entire packet at each hop before transmitting it onward.
    \item[\textit{Cut-Through Routing}] Packet-forwarding method that begins sending to the next hop as soon as enough of the header has arrived to determine the route.
    \item[\textit{Wormhole Flow Control}] Flit-level forwarding method in which one packet may stretch across several routers and links at once.
    \item[\textit{Head-of-Line Blocking}] Throughput loss caused when one blocked packet prevents unrelated traffic behind it from advancing.
    \item[\textit{Virtual Channel}] Logical subdivision of one physical link into multiple independently buffered and arbitrated channels.
    \item[\textit{Router}] Interconnect component that examines incoming traffic, chooses outputs, and manages buffering and arbitration.
    \item[\textit{Data Race}] Incorrect unsynchronized overlap of conflicting accesses to shared state.
    \item[\textit{Embarrassingly Parallel}] Workload whose useful work is effectively all parallelizable and therefore scales nearly ideally if resources permit.
    \item[\textit{Pipelining}] Temporal parallelism where different stages process different instructions simultaneously.
    \item[\textit{Superscalar (Multi-Issue)}] Pipeline organization that can issue more than one instruction per cycle.
    \item[\textit{Register Renaming}] Introduction of additional physical register identities so later writes do not falsely block earlier reads or writes.
    \item[\textit{Op}] One encoded operation within a VLIW instruction stream.
    \item[\textit{Multi-Op}] Group of ops intended to execute in the same cycle in a VLIW design.
    \item[\textit{Tail Bit}] Per-op field marking whether the current op ends the current multi-op.
    \item[\textit{Pause Field}] Encoded delay field that inserts a specified number of idle cycles before the next multi-op begins.
    \item[\textit{Iterative Modulo Scheduling (IMS)}] Practical software-pipelining algorithm that searches for a legal modulo schedule by trying increasing initiation intervals.
    \item[\textit{Initiation Interval (II)}] Number of cycles between the start of successive loop iterations in a software-pipelined schedule.
    \item[\textit{Minimum Initiation Interval (MII)}] Lower bound on II, equal to $\max(ResII,RecII)$.
    \item[\textit{Resource-Constrained Initiation Interval (ResII)}] Lower bound on II imposed by the number of available functional units of each type.
    \item[\textit{Recurrence-Constrained Initiation Interval (RecII)}] Lower bound on II imposed by loop-carried dependence cycles.
    \item[\textit{Prologue}] Non-repeating prefix of a software-pipelined loop that fills the pipeline before steady state.
    \item[\textit{Kernel}] Repeating steady-state portion of a software-pipelined loop.
    \item[\textit{Epilogue}] Non-repeating suffix of a software-pipelined loop that drains remaining work after loop exit.
    \item[\textit{Point Latency Model}] Scheduling model that treats an op as occupying one issue point and becoming ready after a fixed latency.
    \item[\textit{Tomasulo's Algorithm}] Out-of-order scheduling method that combines renaming, reservation stations, and result broadcast to wake dependent instructions when operands become ready.
    \item[\textit{Reservation Station}] Scheduling entry that holds an instruction, its operand state, and dependency tags until it can execute.
    \item[\textit{Common Data Bus (CDB)}] Shared result-broadcast path used to distribute a completed value and its producer tag to waiting instructions.
    \item[\textit{Result Bus}] Generalized form of the CDB; modern designs often use more than one such broadcast path.
    \item[\textit{Architected Register}] Register name defined by the ISA and visible to software.
    \item[\textit{Architected Register File}] Register file containing the committed architectural state visible to software.
    \item[\textit{Physical Register}] Internal hardware register used to hold renamed values beyond the architected register set.
    \item[\textit{Register Alias Table (RAT)}] Mapping table from architected register names to the current physical register locations.
    \item[\textit{Ready Bit}] Status bit indicating that a register or operand value has already been produced and may be consumed.
    \item[\textit{Reorder Buffer (ROB)}] In-order retirement structure that records speculative results and enables precise exceptions.
    \item[\textit{Future File}] Register organization that holds the most recent speculative values while a separate architected register file holds committed state.
    \item[\textit{Messy Register File}] Informal name for the current speculative register state that may contain out-of-order completions not yet reflected in a coherent sequential state.
    \item[\textit{Precise Exception}] Exception model in which all older instructions appear committed and all younger instructions appear not to have changed architectural state.
    \item[\textit{Imprecise Exception}] Exception model that reports an error without guaranteeing an exact sequential recovery point.
    \item[\textit{Checkpoint/Repair}] Recovery method that periodically captures coherent register state and restores it when an exception or mis-speculation occurs.
    \item[\textit{Instruction Boundary}] Dynamic marker used to denote a coherence point or recovery boundary in checkpoint-based schemes.
    \item[\textit{Hyperthreading (SMT)}] Simultaneous multithreading in which multiple threads share one core's execution resources.
    \item[\textit{Speculative Side Channel}] Information leak caused when squashed speculative work still changes microarchitectural state such as caches or predictors.
    \item[\textit{Runtime}] Time to complete one job (response/turnaround time).
    \item[\textit{Throughput}] Completed work per unit time.
    \item[\textit{Amdahl's Law}] Speedup limit for fixed-size workloads when only a fraction is accelerated.
    \item[\textit{Gustafson's Law}] Scaled-speedup model where workload size grows with processor count.
    \item[\textit{Average Access Time (AAT)}] Expected memory-access latency, typically modeled as hit time plus miss rate times miss penalty.
    \item[\textit{SRAM}] Fast, low-latency, higher-area memory used for on-chip caches.
    \item[\textit{DRAM}] Dense, lower-cost, higher-latency memory used for main memory capacity.
    \item[\textit{Hit Time (HT)}] Latency to serve an access from the current cache level.
    \item[\textit{Miss Ratio (MR)}] Fraction of accesses at a level that miss.
    \item[\textit{Miss Penalty (MP)}] Extra latency to service a miss from lower levels.
    \item[\textit{Compulsory Miss}] First-reference miss for data not previously loaded.
    \item[\textit{Capacity Miss}] Miss caused by working set exceeding effective cache capacity.
    \item[\textit{Conflict Miss}] Miss caused by mapping/placement constraints despite sufficient total capacity.
    \item[\textit{Write-Back}] Write policy where modified data is propagated on eviction.
    \item[\textit{Write-Through}] Write policy where lower level is updated on each write hit.
    \item[\textit{Write-Allocate}] Write-miss policy that fetches and installs the block before write.
    \item[\textit{Write-No-Allocate}] Write-miss policy that bypasses install and writes lower level directly.
    \item[\textit{Qubit}] Quantum bit whose useful state must be isolated and controlled well enough to preserve quantum behavior during computation.
    \item[\textit{Decoherence}] Loss of intended quantum behavior when a qubit interacts with the surrounding environment.
    \item[\textit{NISQ}] Noisy intermediate-scale quantum computing; near-term machines with limited qubits, noise, and short usable programs.
    \item[\textit{Logical Qubit}] Error-corrected qubit represented by many physical qubits and maintained through repeated decoding.
    \item[\textit{Phase Oracle}] Reversible quantum operation that encodes a classical function by changing the phase of selected basis states.
    \item[\textit{Fault-Tolerant Quantum Computing}] Quantum-computing approach that uses encoded logical qubits and error correction to run longer computations reliably.
    \item[\textit{Superconducting Logic}] Logic family that uses superconducting devices and very low-resistance wires, often representing information with pulses.
    \item[\textit{Rapid Single Flux Quantum (RSFQ)}] Superconducting pulse-logic family that represents bits using quantized magnetic-flux pulses.
    \item[\textit{Reciprocal Quantum Logic (RQL)}] AC-powered superconducting logic family related to RSFQ and used as one candidate technology for energy-efficient processors.
    \item[\textit{Operand Queue}] Storage structure that holds produced values until compiler-scheduled consumer instructions need them.
    \item[\textit{Barrel Processor}] Multithreaded processor that issues from different resident threads in rotation to keep a deep pipeline occupied.
    \item[\textit{Cryostat}] Cooling system used to maintain cryogenic temperatures for superconducting or quantum hardware.
    \item[\textit{Energy-Delay Product (EDP)}] Efficiency metric combining energy and runtime; lower EDP indicates a better energy/performance tradeoff.
    \item[\textit{SUP}] Superconducting Unit of Performance; proposed organization combining latency-oriented and throughput-oriented superconducting resources.
    \item[\textit{PIPT}] Physically indexed, physically tagged cache organization in which TLB translation typically precedes L1 indexing/tag lookup.
    \item[\textit{VIPT}] Virtually indexed, physically tagged cache organization that overlaps indexing with translation.
    \item[\textit{TLB}] Translation lookaside buffer storing recent virtual-to-physical translations.
    \item[\textit{Structural Hazard}] Stall risk from insufficient hardware resources.
    \item[\textit{Data Hazard}] Stall risk from operand read/write ordering conflicts.
    \item[\textit{Control Hazard}] Stall risk from unresolved branch/jump control flow.
    \item[\textit{RAW}] Read-after-write dependence hazard.
    \item[\textit{WAR}] Write-after-read dependence hazard.
    \item[\textit{WAW}] Write-after-write dependence hazard.
    \item[\textit{Data Dependence Graph}] Graph representation of instruction ordering constraints induced by true and false dependences.
    \item[\textit{Dependence Height}] Priority heuristic measuring the remaining latency on the longest dependent path below an instruction.
    \item[\textit{List Scheduling}] Heuristic compiler scheduling method that repeatedly places the highest-priority ready instructions into the earliest legal cycle.
    \item[\textit{Memory Disambiguation}] Compile-time or runtime reasoning used to determine whether two memory references may access the same location.
    \item[\textit{Forwarding (Bypassing)}] Sending a produced value directly to a consumer before register write-back.
    \item[\textit{Branch Target Buffer (BTB)}] Predictor structure that supplies likely branch targets early in fetch.
    \item[\textit{Aliasing}] Predictor or cache entry interference from unrelated addresses mapping to the same slot.
    \item[\textit{Return-Address Stack (RAS)}] Small stack predictor for call/return target addresses.
    \item[\textit{Direction Prediction}] Branch predictor decision for taken versus not taken.
    \item[\textit{One-Bit Predictor}] Dynamic branch predictor that records only the last observed branch outcome.
    \item[\textit{Smith Predictor}] Saturating-counter branch predictor, commonly implemented with a two-bit counter per entry.
    \item[\textit{GShare}] Global-history predictor that uses an exclusive-or of low PC bits and global history to index the prediction table.
    \item[\textit{GSelect}] Global-history predictor that forms its index by concatenating selected PC bits and global-history bits.
    \item[\textit{Yeh/Patt Predictor}] Two-level branch predictor that uses recorded branch history to index a table of direction counters.
    \item[\textit{Perceptron Predictor}] Branch predictor that computes a weighted sum over branch-history bits and predicts from the sign of that sum.
    \item[\textit{Hybrid Predictor}] Predictor that chooses among multiple underlying predictors, often through a tournament-style selector.
    \item[\textit{BTFNT}] Backward taken, forward not taken static branch heuristic.
    \item[\textit{Delay Slot}] ISA semantic in which the instruction after a branch executes before control transfer takes effect.
    \item[\textit{Squash Slot}] Post-branch instruction position that may be conditionally killed based on branch outcome/policy.
    \item[\textit{WBWA}] Write-back, write-allocate pairing.
    \item[\textit{WTWNA}] Write-through, write-no-allocate pairing.
    \item[\textit{Blocking Cache}] Cache behavior where an outstanding miss prevents servicing younger accesses.
    \item[\textit{Non-Blocking Cache}] Cache behavior allowing selected younger accesses to proceed during a miss.
    \item[\textit{Hit-Under-Miss}] Ability to serve cache hits while a miss is still outstanding.
\end{description}


\section*{Acronym List}\label{sec:acronym-list}
\begin{description}
    \item[AAT] Average Access Time
    \item[ALU] Arithmetic Logic Unit
    \item[AVX] Advanced Vector Extensions
    \item[ACK] Acknowledgment
    \item[AW] Address Width
    \item[BTFNT] Backward Taken, Forward Not Taken
    \item[BTB] Branch Target Buffer
    \item[CDB] Common Data Bus
    \item[CISC] Complex Instruction Set Computer
    \item[CPI] Cycles Per Instruction
    \item[CPU] Central Processing Unit
    \item[CT] Cycle Time
    \item[DDG] Data Dependence Graph
    \item[DSP] Digital Signal Processor
    \item[DISP] Dispatch and Decode
    \item[DRAM] Dynamic Random-Access Memory
    \item[EDP] Energy-Delay Product
    \item[EX] Execute
    \item[FICI] Fire In Order, Complete In Order
    \item[FIFO] First In, First Out
    \item[FICO] Fire In Order, Complete Out of Order
    \item[FOCO] Fire Out of Order, Complete Out of Order
    \item[FP] Floating Point
    \item[FU] Functional Unit
    \item[GPU] Graphics Processing Unit
    \item[HT] Hit Time
    \item[HTM] Hardware Transactional Memory
    \item[IC] Instruction Count
    \item[ID] Instruction Decode
    \item[IF] Instruction Fetch
    \item[II] Initiation Interval
    \item[ILP] Instruction-Level Parallelism
    \item[IMS] Iterative Modulo Scheduling
    \item[IPC] Instructions Per Cycle
    \item[ISA] Instruction Set Architecture
    \item[LFU] Least Frequently Used
    \item[LL/SC] Load-Link / Store-Conditional
    \item[LIP] Least Recently Used Insertion Policy
    \item[LLC] Last-Level Cache
    \item[LRU] Least Recently Used
    \item[MEM] Memory
    \item[MII] Minimum Initiation Interval
    \item[MIMD] Multiple Instruction, Multiple Data
    \item[MIP] Most Recently Used Insertion Policy
    \item[MISD] Multiple Instruction, Single Data
    \item[MESI] Modified, Exclusive, Shared, Invalid
    \item[MP] Miss Penalty
    \item[MR] Miss Ratio
    \item[MRU] Most Recently Used
    \item[MSI] Modified, Shared, Invalid
    \item[MOESI] Modified, Owned, Exclusive, Shared, Invalid
    \item[MOESIF] Modified, Owned, Exclusive, Shared, Invalid, Forward
    \item[MPI] Message Passing Interface
    \item[NACK] Negative Acknowledgment
    \item[NISQ] Noisy Intermediate-Scale Quantum
    \item[NoC] Network on Chip
    \item[OpenMP] Open Multi-Processing
    \item[OS] Operating System
    \item[PC] Program Counter
    \item[PCIe] Peripheral Component Interconnect Express
    \item[PE] Processing Element
    \item[PFN] Physical Frame Number
    \item[PGAS] Partitioned Global Address Space
    \item[PIPT] Physically Indexed, Physically Tagged
    \item[RAS] Return-Address Stack
    \item[RAT] Register Alias Table
    \item[RAW] Read After Write
    \item[RecII] Recurrence-Constrained Initiation Interval
    \item[ResII] Resource-Constrained Initiation Interval
    \item[RISC] Reduced Instruction Set Computer
    \item[ROB] Reorder Buffer
    \item[RMW] Read-Modify-Write
    \item[RQL] Reciprocal Quantum Logic
    \item[RS] Reservation Station
    \item[RSFQ] Rapid Single Flux Quantum
    \item[SIMD] Single Instruction, Multiple Data
    \item[SIMT] Single Instruction, Multiple Threads
    \item[SISD] Single Instruction, Single Data
    \item[SMT] Simultaneous Multithreading
    \item[SRAM] Static Random-Access Memory
    \item[STM] Software Transactional Memory
    \item[SSE] Streaming SIMD Extensions
    \item[SSD] Solid-State Drive
    \item[SU] State Update
    \item[SUP] Superconducting Unit of Performance
    \item[TLP] Thread-Level Parallelism
    \item[TLB] Translation Lookaside Buffer
    \item[TM] Transactional Memory
    \item[VIPT] Virtually Indexed, Physically Tagged
    \item[VLIW] Very Long Instruction Word
    \item[VPN] Virtual Page Number
    \item[WAR] Write After Read
    \item[WAW] Write After Write
    \item[WB] Write Back
    \item[WBWA] Write-Back, Write-Allocate
    \item[WTWNA] Write-Through, Write-No-Allocate
\end{description}
